\textbf{Predstavila: Špela Hajšek}

Naj bo $X$ poljubna neprazna množica. Označimo z $X^{\mathbb{N}}$ množico vseh zaporedij  $(x_{1}, x_{2}, \dots)$ točk iz $X$. Definirajmo razdaljo $d(x,y)$ dveh elementov $x=(x_{1}, x_{2}, \dots)$ in $y=(y_{1}, y_{2}, \dots)$ te množice takole: \[d(x,y) = \begin{cases} 0, & \text{če je } x=y;  \\ \frac{1}{n} & \text{če veljajo } x_{1}=y_{1}, x_{2}= y_{2}, \dots, x_{n-1} = y_{n-1}, \text{ in } x_{n} \neq y_{n}. \end{cases}\]
\begin{enumerate}
	\item Dokažite, da res je $d$ metrika na množici $M$.
	\item Pokažite, da ima ta prostor tole čudno lastnost: vsaka točka poljubne (odprte ali zaprte) krogle je središče te krogle; poleg tega sta vsaki dve krogli bodisi disjunktni ali pa ena od njiju vsebuje drugo.
\end{enumerate}